\chapter{Saisonnalité et effets calendaires}\label{ch:s1:seasonality-and-calendar-effects}

Des centaines d'\omterm{def:s1:seasonality-and-calendar-effects:calendar}{effets calendaires} ont été publiés~: les lundis, les mois de janvier, le passage de fin de mois, les jours précédant les jours fériés, les moitiés de mois, les
saisons. Il suffit d'en chercher assez pour que certains paraissent significatifs dans n'importe quel échantillon. Lakonishok et Smidt, sur quatre-vingt-dix ans de l'indice Dow Jones Industrial
Average, ont trouvé des rendements durablement anormaux autour du passage de la semaine, du mois, de l'année et des jours fériés~; le gaz naturel coûte plus cher en hiver
pour une raison que tout le monde connaît. Sur le marché synthétique de ce chapitre, deux \omterm{def:s1:seasonality-and-calendar-effects:calendar}{effets calendaires} sont insérés parmi cent règles~: le passage de fin de mois résiste
à toutes les corrections pour tests multiples, l'effet précédant les jours fériés est trop faible pour être détecté en trente ans, et la recherche sans correction signale huit
effets qui n'existent pas. La construction s'appelle \texttt{firm.seasonal}.

\begin{definition}[Effet calendaire, saisonnalité des rendements]\label{def:s1:seasonality-and-calendar-effects:calendar}
Un \emph{effet calendaire}\index{effet calendaire} est un écart entre les rendements moyens des jours, des semaines ou des mois définis par le seul calendrier (un
jour de la semaine, un jour du mois, un mois, les jours autour d'un jour férié) et les rendements moyens des autres jours. La \emph{saisonnalité des rendements}%
\index{saisonnalité des rendements} est la récurrence du rendement relatif d'un titre au même moment de chaque année, par exemple le même mois civil.
\end{definition}

\section{Saisonnalité des matières premières pour une raison physique}

Certaines saisons sont physiques. Le gaz naturel est brûlé pour le chauffage en hiver et stocké en été. À partir du prix spot du Henry Hub, de 2000 à 2026 (EIA, via
FRED), l'écart du logarithme du prix par rapport à sa moyenne mobile centrée sur douze mois vaut $+8.8~\%$ en janvier et $+7.3~\%$ en décembre, et $-5.3~\%$ en mars et
$-4.2\%$ en avril, à mesure que les stocks se reconstituent (la \cref{fig:s1:seasonality-and-calendar-effects:months}). Ce schéma n'est pas une stratégie. Les contrats à terme pour une livraison en janvier
sont cotés pour janvier~: la courbe anticipe la saison, comme le font par construction les matières premières saisonnières de \texttt{firm.synthfut}, de sorte qu'une position sur contrats à terme
ne rapporte la saison que si elle était mal cotée. Ce qu'un trader peut détenir, c'est le risque~: un spread hiver--été qui rapporte lorsque l'hiver s'avère plus froid, ou
un stockage plus tendu que la courbe ne l'a anticipé.

Le pétrole brut est le contre-exemple~: sa demande est moins saisonnière et ses stocks sont abondants. En testant le rendement roulant moyen du WTI de chaque mois civil contre celui
des autres mois, de 1986 à 2024, les plus faibles sont octobre ($-2.8~\%$ par mois, $t = -1.71$) et novembre ($-4.2~\%$, $t = -2.17$, une p-valeur brute de 0.03). Douze
mois ont été testés~; après la correction de Bonferroni ou celle de Benjamini et Hochberg, la plus petite p-valeur ajustée vaut 0.36. Il n'y a rien là.

\begin{omfigure}
\begin{tikzpicture}
\begin{axis}[width=10.5cm,height=5cm,ybar,bar width=8pt,xlabel={\small mois civil},ylabel={\small écart à la tendance (\%)},
  tick label style={font=\footnotesize},xtick={1,2,3,4,5,6,7,8,9,10,11,12},xmin=0.4,xmax=12.6,ymin=-7,ymax=11,grid=major,grid style={gray!25},
  table/col sep=comma]
\addplot[fill=omThm!60,draw=omThm] table[x=month,y=gas_pct]{figdata/strategies-1/23-seasonality-and-calendar-effects/months.csv};
\end{axis}
\end{tikzpicture}

\omcaption{Profil saisonnier du prix spot du gaz naturel de Henry Hub~: écart moyen du logarithme du prix mensuel par rapport à sa moyenne mobile centrée sur douze
mois, par mois civil, de janvier 2000 à juillet 2026 (EIA via FRED, série MHHNGSP). Données~: \texttt{s1\_seasonal.gas\_profile}.}\label{fig:s1:seasonality-and-calendar-effects:months}
\end{omfigure}

\section{Effets de fin de mois et effets des jours fériés}

\begin{definition}[Effet de fin de mois]\label{def:s1:seasonality-and-calendar-effects:tom}
L'\emph{effet de fin de mois}\index{effet de fin de mois} est la tendance des rendements boursiers à être plus élevés le dernier jour de négociation du
mois et les premiers jours du mois suivant que les autres jours~; les flux de fin de mois (salaires, cotisations de retraite et rééquilibrages de fonds investis au
début de mois) sont le mécanisme candidat.
\end{definition}

Le test du chapitre est conçu pour rester honnête vis-à-vis de la recherche. \texttt{firm.seasonal} décrit un calendrier stylisé (douze mois de 21 jours de négociation, semaines
de cinq jours, neuf jours fériés par an) sur trente ans de la classe d'actions de \texttt{firm.synthfut} (la moyenne de ses dix contrats à terme sur indices, 14.7~\% de volatilité
annuelle). Deux effets sont insérés~: 8 points de base par jour, le dernier jour et les trois premiers jours de chaque mois, et 15 points de base sur chaque jour précédant un
jour férié. Cent règles sont générées (le \cref{lst:s1:seasonality-and-calendar-effects:rules})~: cinq jours de la semaine, douze mois, 21 jours du mois, douze
fenêtres de fin de mois, les jours avant et après les jours fériés, les semaines du mois, les quinzaines de l'année, les premiers et derniers jours du mois, et quelques
favoris saisonniers («~mai à octobre~», «~décembre et janvier~»). Chaque règle compare le rendement moyen de ses jours à celui des autres. Sur les
cent, 23 recoupent assez les jours de fin de mois pour porter cet effet (une différence attendue d'au moins 3 points de base), 2 portent
l'effet précédant les jours fériés, et 75 ne portent aucun effet.

\begin{center}
\small
\begin{tabular}{@{}l>{\centering\arraybackslash}p{3.6cm}>{\centering\arraybackslash}p{3.6cm}>{\centering\arraybackslash}p{3.6cm}@{}}
\toprule
correction (5~\%) & règles de fin de mois détectées & règles précédant un jour férié détectées & règles sans effet détectées\\
\midrule
aucune & 19 sur 23 & 0 sur 2 & 8 sur 75\\
Bonferroni & 6 & 0 & 0\\
Holm & 6 & 0 & 0\\
Benjamini--Hochberg & 11 & 0 & 3\\
Benjamini--Yekutieli & 6 & 0 & 0\\
\bottomrule
\end{tabular}
\end{center}

Le passage de fin de mois, inséré sur 48 jours par an, est détecté par toutes les corrections~: sa meilleure règle, le dernier jour et les trois premiers jours, a $t = 4.4$
(la \cref{fig:s1:seasonality-and-calendar-effects:rules}). L'effet précédant les jours fériés, sur neuf jours par an, a $t = 1.5$~: réel, plus fort par jour, et invisible,
car neuf jours par an pendant trente ans font 270 observations. Sans correction, la recherche signale huit faux effets, parmi lesquels deux jours de la semaine, trois
jours du mois, une semaine du mois et deux quinzaines, chacun de ces effets figurerait dans le tableau d'un article~; Benjamini--Hochberg, qui contrôle la part des fausses découvertes
plutôt que la probabilité d'en produire au moins une, en retient trois~; les corrections au sens du taux d'erreur familial n'en retiennent aucune. Le livre~4, chapitre~12 a établi les corrections~; ici elles décident
de ce qu'un chercheur croit.

\begin{omfigure}
\begin{tikzpicture}
\begin{axis}[width=11cm,height=5.4cm,xlabel={\small règle, classée par $|t|$},ylabel={\small $|t|$},tick label style={font=\footnotesize},
  xmin=0,xmax=101,ymin=0,ymax=5,grid=major,grid style={gray!25},legend columns=3,legend style={font=\scriptsize,at={(0.5,-0.3)},anchor=north,
  draw=none,fill=none,/tikz/every even column/.append style={column sep=8pt}},unbounded coords=discard,table/col sep=comma]
\addplot[only marks,mark=*,mark size=1.6pt,omThm] table[x=rank,y=tom]{figdata/strategies-1/23-seasonality-and-calendar-effects/rules.csv};
\addplot[only marks,mark=square*,mark size=1.8pt,omDef] table[x=rank,y=pre]{figdata/strategies-1/23-seasonality-and-calendar-effects/rules.csv};
\addplot[only marks,mark=o,mark size=1.4pt,black!60] table[x=rank,y=none]{figdata/strategies-1/23-seasonality-and-calendar-effects/rules.csv};
\addplot[black!50,dashed,forget plot] coordinates {(0,1.96) (101,1.96)};
\addplot[black,dashed,forget plot] coordinates {(0,3.48) (101,3.48)};
\legend{règles de fin de mois,règles précédant les jours fériés,règles sans effet}
\end{axis}
\end{tikzpicture}

\omcaption{Cent règles calendaires testées sur trente ans de la classe d'actions synthétique, avec deux effets insérés~: la valeur absolue de la t-statistique de chaque règle,
classées selon ce qu'elles portent. La ligne en pointillés inférieure est le seuil non corrigé de 5~\% (1.96), la supérieure celui de Bonferroni pour une centaine de tests (3.48). Données~:
\texttt{s1\_seasonal.calendar\_test}.}\label{fig:s1:seasonality-and-calendar-effects:rules}
\end{omfigure}

\section{Saisonnalité au sein d'un même mois dans les rendements des actions}

Heston et Sadka ont mis en évidence une régularité en coupe transversale~: les actions à rendement élevé dans un mois civil ont tendance à avoir un rendement élevé dans le même mois civil
les années suivantes, avec des décalages annuels allant jusqu'à vingt ans, indépendamment de la taille, du secteur, des annonces de résultats, des dividendes et de l'exercice fiscal. Le signal est
simple~: le rendement moyen d'une action dans le même mois civil des années passées. Le chapitre l'insère dans \texttt{firm.synthmkt}~: chaque valeur cotée reçoit un
rendement pour chaque mois civil, identique d'une année à l'autre, avec un écart-type transversal de 1~\% par mois, réparti sur les jours du mois.

Le signal, moyenné sur les cinq années précédentes, présente un IC de rang mensuel de 0.031 avec le rendement du mois ($t = 3.8$ sur 108 mois)~; un portefeuille par déciles
long--short rapporte 0.50~\% par mois, avec un ratio de Sharpe de 0.72, avant coûts. L'effet inséré est petit face à une volatilité spécifique mensuelle de plusieurs
pour cent, et le signal ne moyenne que cinq observations passées de cet effet. Un effet saisonnier réel et stable est un signal faible, car chaque action ne le montre qu'une fois par an.

\section{Ce qui survit hors échantillon}

Sullivan, Timmermann et White ont intitulé leur étude des \omterm{def:s1:seasonality-and-calendar-effects:calendar}{effets calendaires} «~dangers du data mining~», et l'expérience du chapitre en montre le danger en
miniature~: huit fausses découvertes sur cent règles, dont trois même sous un contrôle du taux de fausses découvertes. Ce qui survit possède trois propriétés. Un mécanisme
qui explique pourquoi les jours diffèrent (flux de fin de mois, stockage, chauffage), énoncé avant le test. Un test qui compte chaque règle essayée, y compris celles qui
ont été examinées puis abandonnées. Et suffisamment d'observations~: un effet portant sur neuf jours par an met environ un demi-siècle à atteindre une statistique t de 3.48 à la
taille retenue dans le chapitre. La plupart des \omterm{def:s1:seasonality-and-calendar-effects:calendar}{effets calendaires} publiés échouent à au moins une de ces trois conditions.

\section{Fiches de stratégie}

\begin{strategyfile}{Écart saisonnier du gaz naturel}\label{strat:s1:seasonality-and-calendar-effects:gas}
\sfield{Qui paie, et pourquoi} Les hedgers de la demande hivernale qui paient une prime pour une livraison hivernale au-delà du prix saisonnier attendu.
\sfield{Instruments et plateformes} Contrats à terme sur gaz naturel pour des livraisons d'hiver et d'été.
\sfield{Signal} L'écart hiver--été rapporté aux niveaux de stockage et à son propre historique.
\sfield{Dimensionnement et exécution} Positions de spread, de taille modeste~; reportées avec la strip.
\sfield{Coûts} Modérés~; fourchettes larges sur les échéances lointaines.
\sfield{Comment elle meurt} Hivers froids~: le risque pour lequel elle est rémunérée.
\sfield{Horizon, capacité, infrastructure} Mois~; données de stockage et de météo.
\sfield{Backtest honnête} Prix des contrats à terme, et non profil saisonnier du spot~; données de stockage telles que publiées.
\sfield{Sources} Aucun chiffre de performance vérifié~; profil spot tiré des données de l'EIA.
\end{strategyfile}

\begin{strategyfile}{Effet de fin de mois}\label{strat:s1:seasonality-and-calendar-effects:tom}
\sfield{Qui paie, et pourquoi} Les flux de fin et de début de mois qui achètent des actions à des jours prévisibles.
\sfield{Instruments et plateformes} Contrats à terme sur indices boursiers.
\sfield{Signal} Le calendrier~: le dernier jour de négociation et les trois premiers.
\sfield{Dimensionnement et exécution} Position acheteuse ces jours-là, neutre le reste du temps~; ou inclinaison de l'exposition d'un portefeuille existant.
\sfield{Coûts} Deux allers-retours par mois.
\sfield{Comment elle meurt} Arbitrage d'un schéma publié~; changements dans le calendrier des salaires et des flux des fonds.
\sfield{Horizon, capacité, infrastructure} Jours~; capacité importante sur les contrats à terme sur indices.
\sfield{Backtest honnête} La fenêtre exacte fixée avant le test~; chaque fenêtre essayée est comptée.
\sfield{Sources} Lakonishok et Smidt (1988)~; ce chapitre~: retrouvé par toutes les corrections lorsqu'il est inséré à 8 points de base par jour.
\end{strategyfile}

\begin{strategyfile}{Saisonnalité au sein d'un même mois civil}\label{strat:s1:seasonality-and-calendar-effects:samemonth}
\sfield{Qui paie, et pourquoi} Flux annuels récurrents et événements propres à chaque société.
\sfield{Instruments et plateformes} Actions.
\sfield{Signal} Le rendement moyen d'une action dans le même mois civil des années passées.
\sfield{Dimensionnement et exécution} Déciles ou classement mensuel long--short, neutre aux facteurs.
\sfield{Coûts} Un rééquilibrage complet chaque mois.
\sfield{Comment elle meurt} Encombrement (crowding)~; les coûts d'une rotation mensuelle.
\sfield{Horizon, capacité, infrastructure} Un mois~; de longues séries de rendements.
\sfield{Backtest honnête} Actions radiées incluses~; retards d'années entières uniquement.
\sfield{Sources} Heston et Sadka (2008)~; ce chapitre~: IC 0.031, spread de déciles de 0.50~\% par mois avant coûts.
\end{strategyfile}

\begin{strategyfile}{Effet de veille de jour férié}\label{strat:s1:seasonality-and-calendar-effects:holiday}
\sfield{Qui paie, et pourquoi} Incertain~: l'humeur des marchés et le rachat de positions vendeuses avant une fermeture sont les explications habituelles.
\sfield{Instruments et plateformes} Contrats à terme sur indices boursiers.
\sfield{Signal} La séance qui précède un jour férié de la bourse.
\sfield{Dimensionnement et exécution} Position acheteuse pour la journée.
\sfield{Coûts} Un aller-retour par jour férié.
\sfield{Comment elle meurt} Elle n'a peut-être jamais existé~: peu d'observations par an.
\sfield{Horizon, capacité, infrastructure} Un jour.
\sfield{Backtest honnête} Le calendrier des jours fériés tel qu'il était chaque année~; assez d'années pour que neuf événements par an aient un sens.
\sfield{Sources} Lakonishok et Smidt (1988)~; ce chapitre~: $t = 1.5$ en trente ans lorsqu'il est inséré à 15 points de base.
\end{strategyfile}

\section{Tutoriel~: quelques-uns survivent}\label{tut:s1:seasonality-and-calendar-effects:tut}

\begin{tutorial}
\textbf{Objectif.} Générer cent règles calendaires, les tester sur un marché synthétique comportant deux effets insérés, corriger de quatre façons et compter les découvertes vraies et fausses~;
tester la saisonnalité au sein d'un même mois dans la coupe transversale~; lire les saisons du gaz naturel et du pétrole brut. \textbf{État final~:} le tableau et les deux
figures.
\begin{tutsteps}
\item \textbf{Règles et leur test}.
\omcode{code/firm/seasonal/firm_seasonal.py}{39}{70}{Cent règles calendaires et le test de Welch de chacune.}\label{lst:s1:seasonality-and-calendar-effects:rules}
\item \textbf{L'expérience}~: insertion, test, correction, et classement de chaque découverte selon ce qu'elle porte.
\omcode{code/strategies-1/23-seasonality-and-calendar-effects/python/s1_seasonal.py}{42}{58}{L'expérience calendaire.}\label{lst:s1:seasonality-and-calendar-effects:experiment}
\item \textbf{Exécuter} \texttt{calendar\_test()}, \texttt{same\_month\_test()}, \texttt{wti\_months()} et \texttt{gas\_profile()}, ainsi que \texttt{fig\_seasonal.py}.
\end{tutsteps}
\textbf{Pour aller plus loin.} Exécuter l'expérience sur cinquante graines et compter la fréquence à laquelle chaque correction retrouve l'effet de veille de jour férié~; ajouter une famille de règles par
année d'historique et observer la croissance des fausses découvertes~; insérer un effet de même mois civil qui s'estompe avec le temps.
\end{tutorial}

\section{Construction~: règles calendaires}\label{bld:s1:seasonality-and-calendar-effects:seasonal}

\begin{build}
\bfield{Rôle} Un calendrier stylisé, un générateur de règles, un testeur de règles, des profils saisonniers et le signal de même mois civil.
\bfield{Interface} \texttt{calendar(years, holidays)}, \texttt{rules(cal)}, \texttt{score\_rules(r, rules)}, \texttt{profile(x, period)},
\texttt{same\_month(M, lags)}.
\bfield{Règles} Chaque règle générée est testée et comptée~; le signal de même mois civil n'utilise que des retards d'années entières.
\bfield{Tests d'acceptation} \texttt{code/firm/seasonal/tests/}~: la forme du calendrier et les cent règles distinctes, un jour inséré retrouvé avec la bonne
p-value, les profils et le signal de même mois civil calculés à la main.
\bfield{Pour aller plus loin} Calendriers réels des bourses~; corrections par bootstrap (statistique maximale) pour les règles corrélées~; décomposition saisonnière des courbes de contrats à terme.
\end{build}

\begin{omsources}
\item S.~L.~Heston et R.~Sadka, « Seasonality in the cross-section of stock returns », \emph{Journal of Financial Economics} 87(2), 2008.
\item J.~Lakonishok et S.~Smidt, « Are seasonal anomalies real? A ninety-year perspective », \emph{Review of Financial Studies} 1(4), 1988.
\item R.~Sullivan, A.~Timmermann et H.~White, « Dangers of data mining: the case of calendar effects in stock returns », \emph{Journal of Econometrics}
105(1), 2001.
\item US Energy Information Administration~: Henry Hub natural gas spot price (via FRED) et NYMEX WTI crude oil futures settlements.
\end{omsources}

\section{Exercices}

\begin{exercise}[$\star$]\label{exo:s1:seasonality-and-calendar-effects:1}
Calculer le nombre attendu de fausses découvertes pour cent règles sans effet testées à 5~\%, puis pour 75 règles qui ne portent rien.
\end{exercise}

\begin{exercise}[$\star$]\label{exo:s1:seasonality-and-calendar-effects:2}
Déterminer le seuil bilatéral que fixe la correction de Bonferroni pour cent tests à 5~\%.
\end{exercise}

\begin{exercise}[$\star$]\label{exo:s1:seasonality-and-calendar-effects:3}
Expliquer pourquoi la prime hivernale du prix spot du gaz naturel ne constitue pas une stratégie de trading.
\end{exercise}

\begin{exercise}[$\star\star$]\label{exo:s1:seasonality-and-calendar-effects:4}
La volatilité quotidienne vaut $14.7~\%/\sqrt{252}$. Calculer l'erreur type d'une différence de moyennes pour 1\,440 jours de fin de mois contre 6\,120 autres, et la
statistique t qu'elle donne avec 8 points de base~; faire de même pour 270 jours de veille de férié contre 7\,290 autres et 15 points de base.
\end{exercise}

\begin{exercise}[$\star\star$]\label{exo:s1:seasonality-and-calendar-effects:5}
Calculer le nombre d'années que l'effet de veille de jour férié exigerait pour que sa t-statistique espérée atteigne 3.48.
\end{exercise}

\begin{exercise}[$\star\star$]\label{exo:s1:seasonality-and-calendar-effects:6}
Expliquer pourquoi Benjamini--Hochberg conserve les fausses découvertes que Bonferroni rejette, et dans quels cas cela est acceptable.
\end{exercise}

\begin{exercise}[$\star\star\star$]\label{exo:s1:seasonality-and-calendar-effects:7}
\emph{Programmation.} Modifier la graine de \texttt{simulate\_futures} dans \texttt{calendar\_test} à 8 et relancer. Quels résultats changent, et
lesquels doit-on s'attendre à voir rester inchangés~?
\end{exercise}

\begin{exercise}[$\star\star\star$]\label{exo:s1:seasonality-and-calendar-effects:8}
\emph{Trouver la faille.} «~Les lundis ont été négatifs dans nos données à $t = 2.4$~; nous vendons à découvert des contrats à terme sur indice chaque lundi.~»
\end{exercise}

\section{Problème~: Une poignée survit}

\begin{problem}[{Problème du week-end --- cent règles calendaires}]\label{pb:s1:seasonality-and-calendar-effects:1}
Les marchés synthétiques du chapitre, les données de l'EIA et les archives publiques.

\medskip\noindent\textbf{Partie I --- Saisons.}
\begin{enumerate}
\item Définir l'\omterm{def:s1:seasonality-and-calendar-effects:calendar}{effet calendaire}, l'\omterm{def:s1:seasonality-and-calendar-effects:tom}{effet de fin de mois} et la \omterm{def:s1:seasonality-and-calendar-effects:calendar}{saisonnalité des rendements}.
\item Donner le profil saisonnier du gaz naturel et l'expliquer.
\item Que met en évidence le test du WTI par mois calendaire~?
\item Quels résultats obtiennent Lakonishok et Smidt~?
\end{enumerate}

\medskip\noindent\textbf{Partie II --- L'expérience.}
\begin{enumerate}[resume]
\item Décrire le calendrier, les effets implantés et les cent règles.
\item Comment classe-t-on les règles selon qu'elles portent un effet ou rien~?
\item Donner les découvertes à chaque correction.
\item Pourquoi l'effet de veille de jour férié n'est-il jamais détecté~?
\end{enumerate}

\medskip\noindent\textbf{Partie III --- La coupe transversale.}
\begin{enumerate}[resume]
\item Que trouvent Heston et Sadka~?
\item Décrire l'effet implanté du même mois et le signal.
\item Donner son IC et le spread par déciles.
\item Pourquoi un effet saisonnier stable est-il un signal faible~?
\end{enumerate}

\medskip\noindent\textbf{Partie IV --- Le verdict.}
\begin{enumerate}[resume]
\item Énoncer le \emph{résultat nommé}~: les effets implantés retrouvés et les fausses découvertes à chaque correction.
\item Énumérer les trois caractéristiques d'un \omterm{def:s1:seasonality-and-calendar-effects:calendar}{effet calendaire} qui survit.
\item Quelle correction utiliser pour une recherche sur cent règles, et pourquoi~?
\item Comment faire un backtest honnête d'une stratégie de fin de mois~?
\item À quoi renvoie l'expression «~dangers de la fouille de données~»~?
\item Quelle fiche de stratégie a le mécanisme le plus faible~?
\item Quel lien ce chapitre entretient-il avec le livre~4, chapitre~12~?
\item En une phrase~: quels \omterm{def:s1:seasonality-and-calendar-effects:calendar}{effets calendaires} méritent qu'on y croie~?
\end{enumerate}
\end{problem}

\section{Questions d'entretien}

\begin{interviewq}[$\star$ \iqroles{chercheur}]\label{iq:s1:seasonality-and-calendar-effects:1}
Vous constatez que les vendredis ont été positifs à $t = 2.2$. Que faites-vous ensuite~?
\end{interviewq}

\begin{interviewq}[$\star\star$ \iqroles{chercheur}]\label{iq:s1:seasonality-and-calendar-effects:2}
Quelle est la différence entre contrôler le taux d'erreur familial (FWER) et contrôler le taux de fausses découvertes (FDR)~?
\end{interviewq}

\begin{interviewq}[$\star\star$ \iqroles{trader}]\label{iq:s1:seasonality-and-calendar-effects:3}
Pourquoi les flux de fin de mois peuvent-ils faire bouger les cours des actions, et comment vérifieriez-vous~?
\end{interviewq}

\begin{interviewq}[$\star\star$ \iqroles{développeur}]\label{iq:s1:seasonality-and-calendar-effects:4}
Que doit gérer correctement un calendrier dans un système de backtest~?
\end{interviewq}

\begin{interviewq}[$\star\star$ \iqroles{chercheur}]\label{iq:s1:seasonality-and-calendar-effects:5}
En quoi la saisonnalité du cours au comptant d'une matière première diffère-t-elle de la \omterm{def:s1:seasonality-and-calendar-effects:calendar}{saisonnalité des rendements} de ses contrats à terme~?
\end{interviewq}

\begin{interviewq}[$\star\star\star$ \iqroles{chercheur}]\label{iq:s1:seasonality-and-calendar-effects:6}
Les jours d'une règle représentent une fraction $f$ de $n$ jours, l'effet vaut $\delta$ par jour de règle et la volatilité quotidienne est $\sigma$. Pouvez-vous dériver la t-statistique espérée de la
différence de moyennes et le nombre de jours nécessaires pour qu'elle atteigne $z$~?
\end{interviewq}
